\documentclass[10pt,a4paper]{amsart}
\usepackage{amsmath,amssymb,amsthm,graphicx}
\usepackage{verbatim}
\newenvironment{reference}{\comment}{\endcomment}
%\newenvironment{reference}{}{}
\newenvironment{slogan}{\comment}{\endcomment}
\newenvironment{history}{\comment}{\endcomment}

% For commutative diagrams we use Xy-pic
\usepackage[all]{xy}

% We use 2cell for 2-commutative diagrams.
\xyoption{2cell}
% Theorem environments.
%
\theoremstyle{plain}
\newtheorem{theorem}[subsection]{Theorem}
\newtheorem{proposition}[subsection]{Proposition}
\newtheorem{lemma}[subsection]{Lemma}

\theoremstyle{definition}
\newtheorem{definition}[subsection]{Definition}
\newtheorem{example}[subsection]{Example}
\newtheorem{exercise}[subsection]{Exercise}
\newtheorem{situation}[subsection]{Situation}

\theoremstyle{remark}
\newtheorem{remark}[subsection]{Remark}
\newtheorem{remarks}[subsection]{Remarks}

\numberwithin{equation}{subsection}

% Macros
%
\def\lim{\mathop{\mathrm{lim}}\nolimits}
\def\colim{\mathop{\mathrm{colim}}\nolimits}
\def\Spec{\mathop{\mathrm{Spec}}}
\def\Hom{\mathop{\mathrm{Hom}}\nolimits}
\def\Ext{\mathop{\mathrm{Ext}}\nolimits}
\def\SheafHom{\mathop{\mathcal{H}\!\mathit{om}}\nolimits}
\def\SheafExt{\mathop{\mathcal{E}\!\mathit{xt}}\nolimits}
\def\Sch{\mathit{Sch}}
\def\Mor{\mathop{\mathrm{Mor}}\nolimits}
\def\Ob{\mathop{\mathrm{Ob}}\nolimits}
\def\Sh{\mathop{\mathit{Sh}}\nolimits}
\def\NL{\mathop{N\!L}\nolimits}
\def\CH{\mathop{\mathrm{CH}}\nolimits}
\def\proetale{{pro\text{-}\acute{e}tale}}
\def\etale{{\acute{e}tale}}
\def\QCoh{\mathit{QCoh}}
\def\Ker{\mathop{\mathrm{Ker}}}
\def\Im{\mathop{\mathrm{Im}}}
\def\Coker{\mathop{\mathrm{Coker}}}
\def\Coim{\mathop{\mathrm{Coim}}}

% Boxtimes
%
\DeclareMathSymbol{\boxtimes}{\mathbin}{AMSa}{"02}

%
% Macros for moduli stacks/spaces
%
\def\QCohstack{\mathcal{QC}\!\mathit{oh}}
\def\Cohstack{\mathcal{C}\!\mathit{oh}}
\def\Spacesstack{\mathcal{S}\!\mathit{paces}}
\def\Quotfunctor{\mathrm{Quot}}
\def\Hilbfunctor{\mathrm{Hilb}}
\def\Curvesstack{\mathcal{C}\!\mathit{urves}}
\def\Polarizedstack{\mathcal{P}\!\mathit{olarized}}
\def\Complexesstack{\mathcal{C}\!\mathit{omplexes}}
% \Pic is the operator that assigns to X its picard group, usage \Pic(X)
% \Picardstack_{X/B} denotes the Picard stack of X over B
% \Picardfunctor_{X/B} denotes the Picard functor of X over B
\def\Pic{\mathop{\mathrm{Pic}}\nolimits}
\def\Picardstack{\mathcal{P}\!\mathit{ic}}
\def\Picardfunctor{\mathrm{Pic}}
\def\Deformationcategory{\mathcal{D}\!\mathit{ef}}
\usepackage{cleveref}
\crefname{originalreference}{original reference}{original references}
\Crefname{originalreference}{Original reference}{Original references}
\makeatletter
\newlabel{definition-formal-modification}{{0GDK}{source}{}{sourcepointer-1}{}}
\newlabel{definition-formal-modification@cref}{{[originalreference][0][]0GDK}{[1][0][]source}{}{}{}}
\makeatother
\crefname{originalreference}{original reference}{original references}
\Crefname{originalreference}{Original reference}{Original references}
\makeatletter
\newlabel{lemma-base-change-formal-modification}{{0GDN}{source}{}{sourcepointer-2}{}}
\newlabel{lemma-base-change-formal-modification@cref}{{[originalreference][0][]0GDN}{[1][0][]source}{}{}{}}
\makeatother
\crefname{originalreference}{original reference}{original references}
\Crefname{originalreference}{Original reference}{Original references}
\makeatletter
\newlabel{theorem-dilatations}{{0GDU}{source}{}{sourcepointer-3}{}}
\newlabel{theorem-dilatations@cref}{{[originalreference][0][]0GDU}{[1][0][]source}{}{}{}}
\makeatother
\title{Formal modifications are locally algebraizable}
\author{The Stacks Project Authors}
\date{stacks-project:restricted.tex:local-snapshot-20261012; GNU FDL 1.2 or later}
\begin{document}
\maketitle
\noindent\textit{Editorial note.} All delivered mathematical text is the author's own, copied from the pinned Stacks Project chapter file restricted.tex; no statement or proof step was rewritten, reordered or removed.\par\medskip
\noindent\textit{Editorial note.} The counted proof is the author's complete argument: the counted statement is the author lemma environment labelled lemma-formal-modifications-locally-algebraic; the proof environment immediately following it is its own proof\par\medskip
\noindent\textit{Editorial note.} The label lemma-formal-modifications-locally-algebraic names the counted statement, and the proof environment immediately following it in the source is delivered verbatim as its proof.\par\medskip
\noindent\textit{Editorial note.} The counted statement is carried by the author's own lemma environment in the chapter source.\par\medskip
\noindent\textit{Allowed standard prerequisites (editorial).}\begin{itemize}
\item stability of formal modifications under adic base change (lemma-base-change-formal-modification of the same chapter, delivered as its own card)
\item the dilatation theorem of the same chapter (theorem-dilatations): a formal modification of the completion of a locally Noetherian formal algebraic space along the closed fibre algebraizes to a proper morphism
\end{itemize}\medskip
\ifdefined\hypertarget\hypertarget{sourcepointer-1}{}\fi
\noindent\textit{Source reference.} Original reference 0GDK: sources/stacks-restricted/pinned-original-source.tex, source lines 6275--6275. Prerequisite pointer invoked at 7405; it names a standard fact of the same chapter that the counted argument uses, and every step of the counted proof built on it is retained verbatim in the delivered body.\par\medskip
\ifdefined\hypertarget\hypertarget{sourcepointer-2}{}\fi
\noindent\textit{Source reference.} Original reference 0GDN: sources/stacks-restricted/pinned-original-source.tex, source lines 6366--6366. Prerequisite pointer invoked at 7426; it names a standard fact of the same chapter that the counted argument uses, and every step of the counted proof built on it is retained verbatim in the delivered body.\par\medskip
\ifdefined\hypertarget\hypertarget{sourcepointer-3}{}\fi
\noindent\textit{Source reference.} Original reference 0GDU: sources/stacks-restricted/pinned-original-source.tex, source lines 7376--7376. Prerequisite pointer invoked at 7427; it names a standard fact of the same chapter that the counted argument uses, and every step of the counted proof built on it is retained verbatim in the delivered body.\par\medskip
\begin{lemma}
\label{lemma-formal-modifications-locally-algebraic}
Let $S$ be a scheme. Let $\mathfrak X' \to \mathfrak X$
be a formal modification (Definition \ref{definition-formal-modification})
of locally Noetherian formal algebraic spaces over $S$. Given
\begin{enumerate}
\item any adic Noetherian topological ring $A$,
\item any adic morphism $\text{Spf}(A) \longrightarrow \mathfrak X$
\end{enumerate}
there exists a proper morphism $X \to \Spec(A)$ of algebraic spaces
and an isomorphism
$$
\text{Spf}(A) \times_{\mathfrak X} \mathfrak X'
\longrightarrow
X_{/Z}
$$
over $\text{Spf}(A)$ of the base change of $\mathfrak X$
with the formal completion of $X$ along the ``closed fibre''
$Z = X \times_{\Spec(A)} \text{Spf}(A)_{red}$ of $X$ over $A$.
\end{lemma}

\begin{proof}
The morphism $\text{Spf}(A) \times_{\mathfrak X} \mathfrak X'
\to \text{Spf}(A)$ is a formal modification by
Lemma \ref{lemma-base-change-formal-modification}.
Hence this follows from Theorem \ref{theorem-dilatations}.
\end{proof}

\end{document}
